\section{Application of Program Outcomes}
Table~\ref{tab:po} maps each Program Outcome and Program Specific Outcome addressed by this case study to its concrete implementation. Every mapping demonstrates how foundational engineering knowledge, modern tools, and professional ethics are exercised within the project.

{\small
\begin{longtable}{@{}>{\centering\arraybackslash}p{1.8cm} p{13.2cm}@{}}
    \caption{Mapping of Program Outcomes and Program Specific Outcomes to the Case Study.}\label{tab:po}\\
    \toprule
    \textbf{Outcome} & \textbf{Outcome Description and Concrete Relevance to the Case Study} \\
    \midrule
    \endfirsthead
    \toprule
    \textbf{Outcome} & \textbf{Outcome Description and Concrete Relevance to the Case Study} \\
    \midrule
    \endhead
\textbf{PO1} & \textbf{Engineering Knowledge:} Applies principles of mathematics, computer science, and engineering fundamentals to formalize, model, and analyze the system architecture, mathematical constraints, and foundational algorithms. In this case study (Deep Learning-Based Medical Image Classification for Early Disease Detection), tensor calculus, backpropagation gradients, and gating dynamics are mathematically formalized for Deep Learning-Based Medical Image Classification for Early Disease Detection. \\
    \addlinespace
\textbf{PO2} & \textbf{Problem Analysis:} Systematically identifies, formulates, and investigates complex computational challenges, diagnosing architectural bottlenecks, concurrency anomalies, latency constraints, and integrity failures through rigorous analytical techniques. In this case study (Deep Learning-Based Medical Image Classification for Early Disease Detection), systematic problem analysis resolves vanishing gradients, spatial occlusion, and high-frequency noise in deep representations. \\
    \addlinespace
\textbf{PO3} & \textbf{Design/Development of Solutions:} Designs scalable system architectures, robust database schemas, object hierarchies, or automated pipelines meeting specified operational needs with proper consideration for efficiency, maintainability, and safety. In this case study (Deep Learning-Based Medical Image Classification for Early Disease Detection), a multi-stage deep learning pipeline spanning tensor augmentation, backbone feature extraction, and optimized inference heads is designed. \\
    \addlinespace
\textbf{PO4} & \textbf{Conduct Investigations of Complex Problems:} Uses research-based knowledge and experimental methodologies, including synthesis of empirical benchmarks, stress testing, and profiling, to draw valid technical conclusions. In this case study (Deep Learning-Based Medical Image Classification for Early Disease Detection), rigorous empirical validation using 5-fold cross-validation, ablation studies, and Grad-CAM interpretability provides experimental verification. \\
    \addlinespace
\textbf{PO5} & \textbf{Modern Tool Usage:} Selects and applies state-of-the-art software frameworks, modeling tools, database management engines, profiling utilities, and development environments to deliver high-performance solutions. In this case study (Deep Learning-Based Medical Image Classification for Early Disease Detection), PyTorch, TensorRT, CUDA acceleration, and high-performance computing environments are proficiently utilized. \\
    \addlinespace
\textbf{PO6} & \textbf{The Engineer and Society:} Evaluates contextual consequences of technical deployments on institutional workflows, user accessibility, social well-being, and legal compliance. In this case study (Deep Learning-Based Medical Image Classification for Early Disease Detection), societal impacts on clinical diagnostic speed, energy grid decarbonization, and intelligent transportation safety are prioritized. \\
    \addlinespace
\textbf{PO8} & \textbf{Ethics:} Applies ethical principles and commits to professional ethics, data privacy, user confidentiality, cryptographic integrity, and institutional accountability. In this case study (Deep Learning-Based Medical Image Classification for Early Disease Detection), ethical requirements of clinical model transparency, algorithmic fairness, and patient confidentiality are strictly maintained. \\
    \addlinespace
\textbf{PO12} & \textbf{Life-Long Learning:} Recognizes the need for, and develops the preparation and ability to engage in independent lifelong learning in the broadest context of technological evolution. In this case study (Deep Learning-Based Medical Image Classification for Early Disease Detection), mastery of advanced deep neural architectures equips engineering students with lifelong competencies in artificial intelligence. \\
    \addlinespace
\textbf{PSO1} & \textbf{Program-Specific (AI/ML Core):} Leverages advanced intelligent computational techniques, algorithmic optimization, predictive analytics, and automated decision-making workflows to enhance system capabilities. In this case study (Deep Learning-Based Medical Image Classification for Early Disease Detection), deep convolutional, recurrent, and transformer representations form the core of modern intelligent software systems. \\
    \addlinespace
\textbf{PSO2} & \textbf{Program-Specific (Applied Engineering):} Designs and implements production-grade end-to-end software pipelines integrating intelligent modules with robust backend infrastructure and enterprise services. In this case study (Deep Learning-Based Medical Image Classification for Early Disease Detection), deployment of quantized neural networks onto embedded edge accelerators demonstrates end-to-end applied engineering excellence. \\
    \addlinespace
    \bottomrule
\end{longtable}}